% Pursuing Stacks (cotes 134-2 to 134-8): archive cote 134-7, site page 79 (PDF page 80), typescript page 521.
% typescript. Draft compared: ps6.tex:3929-3995.
% AI-transcribed by Claude Opus 5.5 (claude-opus-5-5, adaptive thinking, effort high), Message Batches API, batch ps4 (msgbatch_012Aoc1NuskATPsyBeKSYK7J), 2026-11-12; not yet checked by eye.
\scanpage{134-7/79}
\mspage{521}

124.   I guess I'll skip giving a\xout{n} more \inserted{(rather!)} or less complete axiomatic
description of $\mathrm{Hot}^{\bullet}_{o}$ or of $\mathrm{Hot}^{\bullet}_{o}(k)$,  in terms of the pair of adjoint
functors $\overset{\hand{\sim}}{\mathrm{LH}}.$ and $\overset{\hand{\sim}}{K}$ and Postnikov dévissage,  although I did go through
this exercise lately - it does'nt really shed new light on the approach
we are following here towards ``schematic'' homotopy types over arbitrary
ground rings.  In contrast with the two other approaches I have been thin%
-king of before (namely De Rham complexes with devided powers,  and small
categories with diagon\xout{x}al maps,  cf. sections 94 and 109 respectively),
the characteristic feature of this approach seems to be that it takes in-%
to account the existence of the canonical functor  \xout{$\mathrm{Hot}_o$xxxD.(Ab)x}
\xout{D.(Ab)xxxHot$_o$}  $\overset{\hand{\sim}}{K} : \mathrm{D.(Ab)} \longrightarrow (\mathrm{Hot})$,  paraphrased by a functor
\[ (1) \qquad \overset{\hand{\sim}}{K} : \mathrm{D.(Ab}_k)(\text{0-conn}) \longrightarrow \mathrm{Hot}_o(k) \]
compatible with ring extension $k \to k'$ ;  while all three approaches have
in common \xout{the existence} that they yield a paraphrase
\[ (2) \qquad \overset{\hand{\sim}}{\mathrm{LH}}. :\ \mathrm{Hot}_o(k) \longrightarrow \mathrm{D.(Ab}_k)(\text{0-conn.}) \]
of the \xout{\unsure{tot}} (total) homology functor,  in a way again compatible with
extension of ground rings.  In the two earlier approaches this \inserted{latter} functor
comes out in a wholly tautological way,\xout{whereas I suspect}  whereas in the
present approach via semisimplicial schemes,  it is'nt quite so tautologi%
cal indeed.  \inserted[typed]{If we want to define a functor $\overset{\hand{\sim}}{K}$ in the context \inserted{say of}} \xout{As a matter of fact in the two earlier approaches, it is}
\inserted[typed]{De Rham algebras with devided \xout{\unsure{induced by}} powers.} \hand{\unsure{k}}
\xout{is easy to define a functor K too as the composition}  one may think
of associating\inserted[typed]{, first to a chain} \xout{D.(Ab$_k$)xxxx\unsure{Hom}} complex in $\mathrm{Ab}_k$  the corresponding semi-%
simplicial k-module,  view this as a semisimplicial set and take it's
De Rham complex with devided powers,  and coefficients in k (as we want
an object over k).  But this is visibly silly,  except possibly for $k=\underline{Z}$,
as the result does'nt depend on the k-module structure of the chain com-%
plex we start with.  It is doubtful anyhow that in this context,  the
tautological functor (2) admits a right adjoint (which we then would call
$\overset{\hand{\sim}}{K}$  of course). - Another noteworthy difference between the \unsure{présent|present} app-%
roach towards defining ``schematic'' homotopy types,  and the two earlier
ones (``earlier'' in these notes, at any rate !)  is that for $k=\underline{Z}$ say,
the category of ``homotopy types over $\underline{Z}$ '' we get here maps into the cate-%
gory of ordinary homotopy types,  whereas it was the opposite with the
two other approaches.

What however I still would like to do,  is a little more daydreaming
about the expected formal properties of the basic functors in between
the two categories appearing in (1) and (2)  - namely,  essentially,
between ``schematic homotopy types''  and ``abelian'' ones.  It will be con-%

\margin{79}
